\section{Conclusion}

This paper introduced the same-input divergence paradigm and applied it to characterize the structure of variation in an AI-driven puzzle hunt generation pipeline. Running the same AoE2 SCOPE.md brief, world knowledge base, and toolkit skill library twice under clean-room conditions produced two complete hunts --- Hunt 1 (``Wololo'') and Hunt 2 (``Run 2'') --- whose outputs we compared at five levels of abstraction spanning round architecture, puzzle type selection, answer word assignment, extraction mechanism design, and individual clue formulations. The resulting empirical characterization contributes four findings to the literature on AI creative generation. First, same-input divergence in this pipeline is structured rather than random: stable and variable decisions partition cleanly into two groups, and that partition is principled rather than incidental. Second, and most importantly, the organizing principle of the partition is not domain knowledge but specification: decisions that are explicitly defined in the creative brief converge across runs; decisions that are left to the pipeline's generative judgment diverge, even when domain constraints narrow the valid option set. This overturns the original domain-constraint hypothesis and points toward specification density as the actionable lever for controlling pipeline predictability. Third, domain semantic gravity is a newly identified phenomenon in which a domain-grounded generator produces thematically coherent word sets across independent runs even when exact string overlap is zero, reflecting a semantic basin of attraction established by the domain's vocabulary and conceptual structure. Fourth, the AI panel's quality signal is moderately stable across the generative decisions that vary most (answer words, mechanisms, clue formulations), confirming that in both observed runs the panel measured structural and cognitive-path quality dimensions that were stable in both runs across the arbitrary choices at Levels 3--5, while remaining sensitive enough to detect genuine quality improvements at the mechanism level (N=2; replication required).

These findings translate directly into practical guidance for practitioners deploying AI creative generators in production pipelines. Decisions that must converge across runs --- scale, format, narrator voice, audience, domain, tone --- should be specified explicitly in the creative brief; in this pipeline and domain, explicitly specified decisions showed stable behaviour across both runs, suggesting a ``run once'' strategy may be appropriate for such decisions. For decisions where variety is desirable (answer word diversity, mechanism type mix), the under-specification should be deliberate, and practitioners should expect and welcome divergence across runs. Unspecified decisions diverged substantially in both runs, consistent with a ``run and select'' strategy for such decisions. For variable decisions that constrain downstream pipeline stages --- most critically, answer words constraining meta mechanism design --- the recommended practice is to run the pipeline multiple times, collect the resulting word sets, and apply human selection among options rather than trusting a single run. These practical recommendations are supported by two runs in a single domain and should be confirmed with N$\geq$5 replication before operational deployment. The same-input divergence paradigm itself is a generalizable method that can be applied to any AI creative pipeline: run the pipeline twice from identical inputs, compare outputs at multiple levels of abstraction, identify which levels show stable behavior (specification-driven) and which show variable behavior (design-space-driven), and use this characterization to inform where in the pipeline human oversight is necessary and where a single AI run is sufficient.

Replication at $N = 10$ runs across three domains --- AoE2, music, and science fiction --- would establish whether the specification-drives-convergence finding generalizes to domains with markedly different knowledge density, vocabulary size, and structural constraints. Extending the comparison framework to include evaluative-language stability --- asking whether the same evaluative frameworks appear in both runs' AI panel reviews, or whether the panel's conceptual vocabulary diverges across runs as the puzzle content does --- would connect the same-input divergence findings to the taxonomy of AI evaluative reasoning developed in the companion profile taxonomy paper \citep{paper3}, and would test whether the stability of panel \textit{scores} is matched by a corresponding stability in panel \textit{reasoning patterns}.
